% =============================================================================
% Reproduction report: attribution graphs on Qwen3.
% Build: `make` (latexmk -pdf). Word count: `make count` (texcount).
%
% WORD BUDGET (limit: 4,000 words total, texcount "words in text"):
%   Abstract                    ~150   (written)
%   1 Introduction              ~640   (written)
%   2 Method                    ~1,450 (written; display math excluded by texcount)
%   3 Experiments and Results   ~1,150 (PLACEHOLDER -- see budget comment in situ)
%   4 Discussion                ~350   (PLACEHOLDER)
%   Slack                       ~250
% Keep the placeholders inside their budgets or trim Sections 1-2 to compensate.
% =============================================================================
\documentclass[11pt,a4paper]{article}

\usepackage[margin=2.4cm]{geometry}
\usepackage[T1]{fontenc}
\usepackage[utf8]{inputenc}
\usepackage{lmodern}
\usepackage{microtype}
\usepackage{amsmath,amssymb,mathtools,bm}
\usepackage{booktabs}
\usepackage{graphicx}
\usepackage[dvipsnames]{xcolor}
\usepackage[numbers,sort&compress]{natbib}
\usepackage[colorlinks=true,linkcolor=BrickRed,citecolor=NavyBlue,urlcolor=NavyBlue]{hyperref}

% ----- notation macros -------------------------------------------------------
\newcommand{\enc}{W_{\mathrm{enc}}}
\newcommand{\dec}{W_{\mathrm{dec}}}
\newcommand{\enci}[2]{W_{\mathrm{enc},#2}^{#1}}   % encoder row #2 at layer #1
\newcommand{\deci}[2]{W_{\mathrm{dec},#2}^{#1}}   % decoder column #2 at layer(s) #1
\newcommand{\bdec}{\bm{b}_{\mathrm{dec}}}
\newcommand{\benc}{\bm{b}_{\mathrm{enc}}}
\newcommand{\jrelu}{\operatorname{JumpReLU}}
\newcommand{\stopgrad}{\bar{J}}          % frozen (stop-gradient) Jacobian
\newcommand{\placeholder}[1]{\begin{quote}\itshape\color{gray!60!black}#1\end{quote}}

\title{Tracing Circuits in Qwen3: A Faithful Reimplementation and Reproduction of Attribution Graphs\thanks{Code: \texttt{llm-circuits} repository. Reference method: \citet{ameisen2025circuit}; reference implementation: \citet{hanna2025circuittracer}.}}
\author{Zixuan Wang\\ \small \texttt{zw499@cam.ac.uk}}
\date{\today}

\begin{document}
\maketitle

% ~150 words -----------------------------------------------------------------
\begin{abstract}
\noindent
Attribution graphs decompose a language model's forward pass on a single prompt into
causal, weighted interactions between sparse dictionary features, and have been used to
map the ``biology'' of a frontier model. We reimplement the full methodology from
scratch for the Qwen3 model family --- local replacement models, batched edge
attribution, influence-based pruning, and constrained-patching interventions --- using
the open-source \mbox{circuit-tracer} library only to load pretrained per-layer
transcoders. The reimplementation is verified numerically equivalent to
circuit-tracer: identical node sets, matching edge weights, identical pruned graphs,
and matching constrained-steering effects on shared inputs. We then reproduce two case studies
from \citet{lindsey2025biology} on Qwen3-4B --- two-digit addition and multilingual
antonyms --- following the paper's intervention protocols, and add a same-recipe
cross-scale comparison (4B vs.\ 8B).
% TODO(results): sync this sentence with Section 4 once it is written.
The core circuit structures reproduce; several interventions reproduce only at reduced
steering strength; and one causal claim (the language swap) localises differently in
Qwen3, which we characterise mechanistically.
\end{abstract}

% =============================================================================
\section{Introduction}\label{sec:intro}
% ~640 words. Related work is folded in here (no separate section).
% =============================================================================

Mechanistic interpretability aims to reverse-engineer the computations of neural
networks into human-understandable algorithms~\citep{olah2020zoom,elhage2021mathematical}.
For large language models a central obstacle is that the natural unit of computation is
not the neuron: models appear to represent many more concepts than they have neurons,
storing them in \emph{superposition}, so individual neurons are polysemantic and
resist interpretation~\citep{elhage2022superposition}. Explaining a behaviour
mechanistically therefore requires both a better feature basis and a way to trace, on
concrete prompts, how those features causally interact.

Dictionary learning addresses the first requirement. Sparse autoencoders (SAEs)
decompose activations into an overcomplete basis of sparsely-firing, substantially more
monosemantic features~\citep{bricken2023monosemanticity,cunningham2024sparse}, scale to
frontier models~\citep{templeton2024scaling}, and benefit from discontinuous activation
functions such as JumpReLU~\citep{rajamanoharan2024jumprelu}. An SAE, however, only
re-describes the representation at one point in the network: it says what is encoded,
not how it is computed. \emph{Transcoders}~\citep{dunefsky2024transcoders} instead
train the same sparse dictionary architecture to imitate an MLP's input--output map, so
that each feature has both an input action (its encoder row) and an output action (its
decoder column) and the MLP's computation factors through an interpretable sparse
bottleneck. Cross-layer variants let one feature write to all subsequent
layers~\citep{ameisen2025circuit}.

The second requirement --- tracing interactions --- has traditionally been met with
causal patching. Manual activation patching over attention heads and MLPs isolated,
e.g., the indirect-object-identification circuit in GPT-2~\citep{wang2023ioi};
ACDC automated the search~\citep{conmy2023acdc}; attribution patching replaced
interventions with cheap gradient approximations~\citep{syed2024attribution}; and
sparse feature circuits lifted patching from model components to SAE
features~\citep{marks2025sparse}. \emph{Attribution
graphs}~\citep{ameisen2025circuit} sharpen this line of work into a per-prompt causal
graph over transcoder features. Replacing MLPs with transcoders, freezing attention
patterns and normalisation denominators at their on-prompt values, and inserting error
nodes yields a \emph{local replacement model} that reproduces the underlying model's
forward pass exactly while being linear in the feature activations; \emph{direct}
effects between features are then exact linear coefficients obtained by
backpropagation rather than approximations, and the resulting dense graph is pruned to
an interpretable subgraph whose hypotheses are validated by steering interventions on
the real model. The companion paper applies the method across dozens of behaviours of
Claude~3.5 Haiku~\citep{lindsey2025biology}, and the circuit-tracer
library~\citep{hanna2025circuittracer} open-sources the pipeline for open-weight
models. Together these form the reference for the present work.

This report makes two contributions. \textbf{First}, we reimplement the entire
pipeline from scratch for HuggingFace Qwen3 models~\citep{yang2025qwen3}: capture of
frozen constants, the local replacement forward pass, batched edge attribution via
\texttt{autograd}, influence-based graph pruning, an interactive graph explorer, and
constrained-patching feature interventions. circuit-tracer is used \emph{only} to load
pretrained per-layer transcoders~\citep{hanna2025qwen3tc}; no attribution, pruning, or
intervention machinery is imported. A dedicated verification suite establishes
numerical equivalence with circuit-tracer on shared inputs --- identical node sets,
edge weights matching to floating-point noise, identical influence rankings and pruned
graphs, and matching constrained-steering deltas (Section~\ref{sec:experiments}). \textbf{Second},
we reproduce two case studies of \citet{lindsey2025biology} on Qwen3-4B: two-digit
\emph{addition} (parallel lookup-table and sum pathways, the operand-plot feature
taxonomy, causal input suppressions, and cross-context reuse of arithmetic features)
and \emph{multilingual antonyms} (a shared cross-lingual ``say-large'' core with
language-specific output stages, operation/operand/language swap interventions, and
cross-lingual feature overlap as a function of scale, compared against Qwen3-8B with
same-recipe transcoders). Where Qwen3's tokenisation or capacity forces protocol
adaptations --- digit-by-digit arithmetic, chat templates, steering strengths at which
a 4B model degrades --- we document the adaptation explicitly and separate model
differences from method failures.

Section~\ref{sec:method} presents the method; Section~\ref{sec:experiments} reports
the verification and the two reproductions; Section~\ref{sec:discussion} discusses
what transfers across models and what does not.

% =============================================================================
\section{Method}\label{sec:method}
% ~1,450 words. Describes the REFERENCE method (Ameisen et al. 2025 +
% circuit-tracer v0.3.1), which our reimplementation follows and is verified
% against. Only: transcoders, attribution-graph generation, pruning,
% interventions.
% =============================================================================

Throughout, the underlying model is a pre-norm decoder-only transformer with $L$
layers. We write $\bm{x}^{\ell}_{c}\in\mathbb{R}^{d}$ for the residual stream at token
position $c$ before layer $\ell$'s MLP, $\tilde{\bm{x}}^{\ell}_{c}$ for the normalised
MLP input,\footnote{In Qwen3, RMSNorm: $\tilde{\bm{x}} = \bm{x} \oslash
\sigma(\bm{x})\odot\bm{\gamma}$ with $\sigma(\bm{x})=\smash{\sqrt{\lVert \bm{x}
\rVert^2/d+\varepsilon}}$.} and $\bm{y}^{\ell}_{c}=\mathrm{MLP}^{\ell}(\tilde{\bm{x}}^{\ell}_{c})$
for the MLP output added back into the stream. Logits are
$\mathrm{logit}_{t} = \bm{u}_{t}^{\top}\,\mathrm{norm}(\bm{x}^{L+1})$ with unembedding
rows $\bm{u}_{t}$.

\subsection{Transcoders: sparse dictionaries for MLP computation}\label{sec:method:tc}

An SAE learns an overcomplete sparse code
$\bm{f}(\bm{x})=\sigma(\enc\bm{x}+\benc)$ with linear decoder
$\hat{\bm{x}}=\dec\bm{f}+\bdec$, trained to reconstruct its own input under a sparsity
penalty. A \emph{transcoder} keeps this architecture but changes the target: it is
trained to predict the MLP's \emph{output} from the MLP's \emph{input}, so that the
learned features mediate the layer's computation rather than merely re-encode its
input. In the per-layer form (PLT) used in this work, layer $\ell$ has feature
activations and reconstruction
\begin{equation}
\bm{a}^{\ell}_{c} \;=\; \jrelu\!\big(\enc^{\ell}\,\tilde{\bm{x}}^{\ell}_{c}+\benc^{\ell}\big),
\qquad
\hat{\bm{y}}^{\ell}_{c} \;=\; {\dec^{\ell}}^{\top} \bm{a}^{\ell}_{c}+\bdec^{\ell},
\qquad
\jrelu(z)_i \;=\; z_i\,\mathbb{1}[z_i>\theta_i],
\label{eq:plt}
\end{equation}
i.e.\ a linear map on either side of a hard sparse gate. The reference method's
default is the \emph{cross-layer} transcoder (CLT), in which a feature encoded at
layer $\ell'$ writes to the MLP outputs of \emph{all} later layers through separate
decoders, $\hat{\bm{y}}^{\ell}_{c}=\sum_{\ell'\le\ell} {\dec^{\ell'\to\ell}}^{\top}
\bm{a}^{\ell'}_{c}+\bdec^{\ell}$; the PLT is the special case
$\dec^{\ell'\to\ell}=0$ for $\ell'\neq\ell$. Training minimises reconstruction error
with a tanh sparsity penalty,
\begin{equation}
\mathcal{L} \;=\; \sum_{\ell=1}^{L}\big\lVert \hat{\bm{y}}^{\ell}-\bm{y}^{\ell} \big\rVert_2^2
\;+\; \lambda \sum_{\ell,i} \tanh\!\big(c\,\lVert \deci{\ell}{i}\rVert\, a^{\ell}_{i}\big),
\end{equation}
with the JumpReLU thresholds $\theta$ trained through a straight-through estimator.
The pretrained Qwen3 transcoders used here are per-layer JumpReLU dictionaries with
$163{,}840$ features per layer and no skip connection~\citep{hanna2025qwen3tc}.

\subsection{Attribution-graph generation}\label{sec:method:attr}

\paragraph{Local replacement model.}
Substituting every MLP by its transcoder yields the \emph{replacement model}, which
only approximates the underlying model. The attribution graph is instead built on the
\emph{local replacement model} for a specific prompt $p$: (i)~MLPs are replaced by
transcoders; (ii)~attention patterns and normalisation denominators are \emph{frozen}
at the values they take in the underlying model's forward pass on $p$, making
attention and normalisation linear maps of the residual stream; and (iii)~an
\emph{error node} is added at every (position, layer),
\begin{equation}
\bm{e}^{\ell}_{c} \;=\; \bm{y}^{\ell}_{c} - \hat{\bm{y}}^{\ell}_{c},
\end{equation}
injected as a constant so that every activation and logit of the local replacement
model \emph{exactly} equals the underlying model's. By construction the only remaining
nonlinearities are the JumpReLU gates on feature pre-activations: between any feature's
(constant-scaled) output and any downstream pre-activation, the map is affine.

\paragraph{Nodes.}
The graph has four node types: one \emph{embedding} node per prompt token; one
\emph{feature} node per active feature instance $s=(c_s,\ell_s,i_s)$ with activation
$a_s>0$; one \emph{error} node per (position, layer); and \emph{logit} nodes for the
smallest set of candidate output tokens whose cumulative probability reaches $0.95$,
capped at $10$, at the final position. Embedding and error nodes have no incoming
edges; logit nodes have no outgoing edges.

\paragraph{Edges.}
Edge weights are the \emph{direct} linear effects of source on target within the local
replacement model. Let $\stopgrad_{(c_s,\ell)\to(c_t,\ell_t)}$ denote the Jacobian of
the residual stream at $(c_t,\ell_t)$ with respect to the stream at $(c_s,\ell)$
computed with stop-gradients on all MLP/transcoder outputs, attention patterns, and
normalisation denominators --- i.e.\ the sum over all paths through the residual
stream and the frozen attention OV circuits, but through no other feature's
nonlinearity. For a feature source $s$ and a feature target $t$ with encoder row
$\enci{\ell_t}{i_t}$, the edge weight is
\begin{equation}
A_{s\to t} \;=\; a_s \sum_{\ell_s \le \ell < \ell_t}
\big(\deci{\ell_s\to\ell}{i_s}\big)^{\top}\,
\stopgrad_{(c_s,\ell)\to(c_t,\ell_t)}\;
\enci{\ell_t}{i_t},
\label{eq:edge}
\end{equation}
which for the PLT collapses to the single $\ell=\ell_s$ term. Embedding and error
sources use their output vectors ($\bm{E}_{\mathrm{tok}}$ and $\bm{e}^{\ell}_c$) in
place of $a_s\deci{}{i_s}$. For a logit target the input vector is the
\emph{vocabulary-demeaned} unembedding direction $\bm{u}_t-\bar{\bm{u}}$, applied at
the (frozen-)normalised final residual --- the demeaned logit is the component that
affects the softmax, which is invariant to uniform shifts. Because the local
replacement model is affine in the source outputs, incoming edges are \emph{complete}:
the pre-activation of any feature node satisfies $h_t=\sum_{s}A_{s\to t}+\text{const}$,
where the constant collects bias-path terms ($\benc$, $\bdec$) that are not
represented as nodes. Attention QK circuits are deliberately outside the graph: nodes
influence one another through frozen attention outputs, but effects \emph{on} the
patterns themselves are not modelled.

\paragraph{Computation.}
Edges are computed target-by-target with one backward pass per batch of targets: the
target's input vector is injected as the gradient seed at its (layer, position), the
backward pass runs through the frozen-linear graph, and every source's edge weight is
the inner product of the resulting gradient at the source's write location with the
source's output vector, as in Eq.~\eqref{eq:edge}. The cost is linear in the number of
attributed targets and independent of the number of sources; on dense prompts, targets
can be selected adaptively in decreasing order of an incrementally re-estimated
influence on the logits, so that the most explanatory feature nodes are attributed
first.

\subsection{Graph pruning}\label{sec:method:prune}

Raw graphs contain $10^4$--$10^6$ edges, most irrelevant to the output. Pruning keeps
the subgraph that carries the bulk of the influence on the logits. Let $A$ be indexed
$[t,s]$ (rows = targets) and define the normalised unsigned adjacency
\begin{equation}
\hat{A}_{ts} \;=\; \frac{\lvert A_{ts}\rvert}{\sum_{s'} \lvert A_{ts'}\rvert},
\qquad
B \;=\; \sum_{k\ge1} \hat{A}^{k} \;=\; (I-\hat{A})^{-1} - I,
\end{equation}
where the Neumann series converges because the graph is a DAG ($\hat A$ is nilpotent);
$B_{tv}$ sums the strengths of all directed paths $v\!\to\!t$, a path's strength being
the product of its normalised edge weights. Each node $v$ receives the
probability-weighted logit influence
\begin{equation}
I_v \;=\; \sum_{t\,\in\,\mathrm{logits}} p_t\, B_{tv},
\label{eq:influence}
\end{equation}
with $p_t$ the model's output probability for logit node $t$. \textbf{Node pruning}
keeps the minimal set of nodes, in decreasing $I_v$, whose cumulative share of
$\sum_v I_v$ reaches $\tau_{\mathrm{node}}=0.8$; embedding and logit nodes are always
kept.\footnote{The paper's prose also protects error nodes; the reference
implementation (and ours, verified equivalent to it) lets error nodes be pruned.}
\textbf{Edge pruning} recomputes $\hat{A}'$ and $I'$ on the pruned graph and scores
each surviving edge by the influence it feeds into its target,
\begin{equation}
s_{ts} \;=\; \hat{A}'_{ts}\,\big(I'_t + p_t\big),
\end{equation}
($p_t=0$ for non-logit targets, and $I'_t=0$ for logit nodes, so the parenthesis is
the target's score in either case), again keeping the minimal set with cumulative
share $\tau_{\mathrm{edge}}=0.98$. Finally, internal nodes left dangling --- features
without both an incoming and an outgoing kept edge, error nodes without an outgoing
one --- are removed iteratively. Typical graphs shrink by roughly an order of
magnitude in nodes and two in edges while retaining $\approx\!80\%$ of the logit
influence.

\subsection{Interventions}\label{sec:method:interv}

Attribution graphs are hypotheses; they are validated by steering features in the
\emph{underlying} model and checking that downstream features and logits move as the
graph predicts. Steering feature $s$ multiplicatively by $M$ --- $a_s \mapsto M a_s$,
so $M{=}0$ ablates and $M{=}-1$ flips the sign\footnote{Our code parameterises the
same operation by the additive multiple $m = M-1$ of the clean activation
($m{=}{-}1$ ablate, $m{=}{-}2$ flip); circuit-tracer's API takes the absolute target
value $Ma_s$. All strengths in this report are quoted in the paper's multiplicative
$M$ convention.} --- adds the corresponding decoder delta to every MLP output the
feature writes to:
\begin{equation}
\Delta \bm{y}^{\ell}_{c_s} \;=\; (M-1)\,a_s\, \deci{\ell_s\to\ell}{i_s},
\qquad \ell_s \le \ell \le \ell_e .
\end{equation}
Under \textbf{constrained patching} over a layer range $[\ell_s,\ell_e]$, all MLP
outputs within the range are pinned to their clean recorded values before the deltas
are added, so the intervention itself is the \emph{only} change to any MLP output in
the range (no second-order effects within it, matching the graph's direct-effect
semantics); computation resumes normally after $\ell_e$. Attention patterns are frozen
to their clean values throughout, consistent with the graph's frozen-attention
assumption; normalisation denominators are frozen only when the range spans the whole
model, in which case the intervention measures a pure direct (linear) effect on the
logits. The alternative, \textbf{iterative patching}, lets downstream features
re-encode their perturbed inputs and take new values, propagating knock-on effects
through the real MLPs at the price of compounding reconstruction errors. Effects are
read out both as changes in the output distribution and as downstream feature
activations re-encoded from the perturbed forward pass, reported as percentages of
their clean baselines. Because feature dictionaries are imperfect (reconstruction
error, incomplete feature recruitment), steering at $\lvert M\rvert$ beyond the
observed activation is often required before behavioural effects appear; the paper
typically suppresses source features to $M=-1$ or beyond and drives donor features at
several times their donor-prompt activation.

\paragraph{Reimplementation.}
Our implementation realises the above for HuggingFace Qwen3 checkpoints with per-layer
transcoders: a capture pass records attention patterns, RMSNorm denominators and MLP
outputs; hooks then replace MLPs with transcoder reconstructions plus constant error
nodes, substitute frozen-pattern attention (V/O only) and frozen-denominator norms;
batched \texttt{autograd} computes Eq.~\eqref{eq:edge} for hundreds of targets per
backward pass. Pruning follows Section~\ref{sec:method:prune} exactly, and
interventions follow circuit-tracer's constrained semantics. Numerical equivalence
with circuit-tracer is established in Section~\ref{sec:experiments}.

% =============================================================================
\section{Experiments and Results}\label{sec:experiments}
% =============================================================================
% PLACEHOLDER -- NOT YET WRITTEN. Budget: ~1,150 words + 2--3 figures/tables.
% Planned narrative (source material: notebooks/addition.ipynb,
% notebooks/multilingual.ipynb, DEVLOG.md, notes/biology_digest.md):
%
%   3.1 Setup (~150 w): Qwen3-4B (36 layers, bf16) + mwhanna per-layer JumpReLU
%       transcoders (163,840 features/layer); A100-80GB; ~8,000-target graph in
%       ~1-2 min; two notebooks re-executed end-to-end with bf16 parity across
%       three A100 variants.
%
%   3.2 Faithfulness verification (~150 w): 11/11 checks vs circuit-tracer on
%       Qwen3-0.6B (forward exactness 1.2e-5; linearity identity 1.7e-5;
%       200 edge spot-checks vs plain autograd 3.7e-5; salient-logit parity;
%       node-set Jaccard 1.0; edge cosine/Pearson 1.0 over 2.4M shared pairs;
%       influence Spearman 1.0; pruned-set Jaccard 1.0) + steering parity
%       (max |dOurs - dCT| = 6e-5, noise floor). One table.
%
%   3.3 Addition (~450 w): digit-split adaptation (Qwen3 tokenises digits
%       individually -> two per-digit circuits); studied pair 46+49=95 (paper's
%       36+59 answered wrong by the 4B model; digit-class-preserving switch);
%       operand-grid taxonomy reproducing the paper's classes (sum
%       anti-diagonals, lookup-table lattices, mod-10, magnitude bands,
%       operand crosses); input-suppression strength ladder (ablation /
%       paper -1x / paper figure -2x) -- phenomenology at reduced strength,
%       two-band magnitude supernode giving the clean ones/magnitude
%       dissociation at full paper strength; lookup-vs-sum steering; polymer
%       citation reuse (8/10 features at the ones moment) + verified-donor
%       lookup swap (8 @ 0.715, paper 998 @ 66.6%); lookups-only suppression
%       silencing all sum features (indirect-effect architecture); computed-9
%       negative-weight screen (measured negative); introspection sampling
%       (carry narration in ~half of samples = paper's mismatch motif).
%
%   3.4 Multilingual (~400 w): shared say-big core across en/fr/zh (107
%       features in all three pruned graphs; trio L30/L31/L32); operation swap
%       (top-1 in all three languages at 1-2x in the paper's raw format;
%       over-drive at full +/-5-6x -- capacity difference); operand swap
%       (cold/f/cold-zh at crossovers 0.625x/1.0x/0.875x -- cleanest
%       reproduction); language swap null on raw prompts with populated early
%       supernodes + Fig-B5-style say-large readouts (say-large-zh never
%       activates; causal handle sits in late say-X features -- an earned
%       model difference); overlap-by-layer IoU with unrelated-pair baseline
%       (paper-shaped, en-fr strongest); same-recipe scale check: 8b > 4b on
%       every language pair (paper's scale claim reproduced once the
%       transcoder-recipe confound is removed).
% =============================================================================
\placeholder{[Placeholder --- to be written; budgeted at $\approx$1,150 words.
Verification of the reimplementation against circuit-tracer, then the addition and
multilingual reproductions from \texttt{notebooks/}, with per-experiment verdicts
against \citet{lindsey2025biology}.]}

% =============================================================================
\section{Discussion}\label{sec:discussion}
% =============================================================================
% PLACEHOLDER -- NOT YET WRITTEN. Budget: ~350 words.
% Planned content: verdict synthesis (what reproduced / at reduced strength /
% differs); model-vs-method attribution of the differences (4B capacity ->
% over-drive at paper strengths; EN/ZH-trained Qwen3 -> shared mechanistic
% default; late say-X features carry output language); limitations (PLT vs
% CLT, one model family, two case studies, chat-format sensitivity, feature
% dictionaries' unexplained variance); what the faithfulness verification
% does and does not license; future work (CLT edges, more behaviours,
% global weights analysis).
% =============================================================================
\placeholder{[Placeholder --- to be written; budgeted at $\approx$350 words.]}

\bibliographystyle{unsrtnat}
\bibliography{refs}

\end{document}
